\section{System}

\subsection{The assembly surface}

We treat the context window as a compilation target, not a transcript buffer.
Every turn the runtime assembles an ordered set of typed sections and hashes the
result before the model sees it. The code does not decide what belongs in the
window; a signed data file does.

That file declares fifteen kinds of section: charter, task, intent, tool
definitions and results, skills, project rules, repository map, locator card,
snippets, history summaries, memory, handoff material, progress document and
intent history. Each kind declares its source, loading mode, trust level, share
of the turn budget, drop policy, verification, the largest unit it may inline,
pointer preview, and the tool that expands a pointer.

Two properties do the work. First, the budget has three levels: a per-turn
total, a share per kind, a cap per unit. Exceeding the unit cap does not
truncate text; it replaces the text with a pointer into an object store, and
the pointer must resolve. Exceeding the share drops the section and writes the
drop to the ledger. A task profile may re-weight shares, but they must still
sum to at most the total; the loader refuses the registry otherwise.

Second, the loading mode carries the argument. Resident kinds --- charter, task,
intent --- are always present. Search and lazy kinds enter only when a query
resolves and hits. Summary kinds carry references back to the original. Derived
kinds are recomputed at assembly time from this turn's evidence and compared
with the recorded derivation, so a card cannot be inherited from an earlier
turn and passed off as fresh.

Before the writer's first token, a preflight step runs the failing selector with
no model call and records the frames, the failing test function and the failure
excerpt, leaving the tree unchanged. Its purpose is not speed: the first
evidence in the window is produced and recorded by the runtime, so the model
cannot be shown a tool result the ledger lacks.

\subsection{Invariants, and what they stop}

The register names twenty-one invariants over the assembly, evaluated at the
entry and fail-closed: a violation rejects the turn rather than degrading it.
Five carry the delivery argument.

The typed-section invariant requires every section to have a declared kind and rejects an
unknown one at the entry. It stops a new category of text --- a pasted log, a
stack trace --- from entering the window because some code path found it
convenient. Without it the share budget is arithmetic over an open set: correct
as arithmetic, hollow as a guarantee.

The query-justification invariant requires inclusion to be justified by a resolvable query and
a match. It stops ``just in case'' context, and it makes inclusion falsifiable:
every section has an answer to the question of why it is there.

The pointer-substitution invariant turns a unit that exceeds its cap into a pointer with a
resolvable reference inside the object store. It stops one file from owning the
window, and it made the ablation legible: a runaway read could no longer
quietly consume the evidence budget, because eviction became a logged drop.

The drop-logging invariant requires a dropped section to be logged. It stops silent
truncation, the failure that would make the arm label false: a context-on arm
can decay into ``context on for the first few turns'' while still being called
context-on.

The payload-hash invariant requires the hash recorded for a turn to be the hash of the
payload the model actually received. It stops post-hoc reconstruction, and it
later exposed a delivery bug: the locator card was written and scheduled, but
never reached the request body (Section~\ref{sec:locator}).

These five are not enforced with equal strength: some are checked structurally,
against the shape of a derivation, some semantically, against whether its
entries hold. We report the difference rather than rounding it away, because an
invariant documented as stronger than it is behaves like no invariant at all.

\subsection{Task classes}

The first runtime knew one kind of work: fix this bug, run the tests, accept or
reject. A task class is now a row in a data file with three fields. The plan
gives the step order and the prompts. The profile gives the budget shares;
long-horizon work moves budget from the repository map and the skill
descriptions into tool results, because there the bottleneck is remembering
what was read, not learning what tools exist. The gate is looked up
in a small registry: tests pass on the model's own report, node-level pass over
an externally supplied selector, or a command whose exit status decides and
whose output summary enters the evidence.

The point is not configurability: the SWE-bench runs are graded by node, with
the selector taken from the dataset. A new class costs a row, not a branch. The
loader refuses unknown fields, plans, profiles and gates, and an absent table
falls back to old behavior, so the abstraction cannot quietly change the
control arm. Our tool-and-policy experiments reuse the assembly surface through
a separate adapter that does not run the plan engine; we state that boundary
rather than paper over it.

\subsection{The accounting surface}

Every assembled turn has a canonical form and a hash; the rendered body is kept
in an object store, and the ledger keeps the hash and the reference, not the
text. The ledger therefore records what was delivered, not what was intended. A
turn that was assembled but never delivered appears as a gap, not a
plausible-looking row.

The execution layer adds epochs. The table of tool plates --- which tools exist,
with declared effect and risk --- is hashed into an epoch; reports carry the
epoch, and counts are never summed across epochs. Changing one line of that
table yields a new epoch, which guards against a measurement improving because
the thing measured changed underneath it.

For the experiments the consequence is narrow. The two arms differ by a single
switch; preflight, the gate, the close-out rule and the accounting are common
to both. The accounting surface is what lets us assert that without relying on
our memory of the runs.

% FACTS-USED: E9 (locator card never reached the request body)
